% lib/ext-eprint-dbx.tex
% Emits the custom biblatex data model for extended eprints (e.g., URN, MR, etc.)
\begin{filecontents*}{ext-eprint.dbx}
% Source - https://tex.stackexchange.com/a/163787
% Posted by moewe, modified by community. See post 'Timeline' for change history
% Retrieved 2026-06-11, License - CC BY-SA 3.0
% 2016/09/11 extended stand-alone eprint fields
% Augmented by GQMJr to add additional fields
\ProvidesFile{ext-eprint.dbx}[2026/08/28 v2.0 Extended persistent identifier data model]

% Use verbatim for URLs, file paths, raw hashes, DOIs, or identifier tokens containing characters that would otherwise conflict with TeX syntax.
% Raw verbatim identifiers
\DeclareDatamodelFields[type=field,datatype=verbatim]{%
  % Scholarly indexing & Repositories
  arxiv,
  mr,
  zbl,
  jstor,
  hdl,
  pubmed,
  pmcid,
  googlebooks,
  adsbibcode,
  osti,
  ntis,
  eric,
  ssrn,
  % National authority & Global library identifiers
  ark,
  urn,
  lccn,
  oclc,
  libris,
  gnd,
  dnb,
  sudoc,
  sbn,
  asin,
  % Proprietary metrics (optional)
  scopus,
  wos%
}

% Literal strings
\DeclareDatamodelFields[type=field,datatype=literal]{%
  arxivclass,
  oai%
}

% URI fields
\DeclareDatamodelFields[type=field,datatype=uri]{%
  archiveurl,
  adsurl%
}

% Date fields
\DeclareDatamodelFields[type=field,datatype=date,skipout]{%
  archivedate%
}

% Register all fields as valid across all entry types
\DeclareDatamodelEntryfields{%
  arxiv,
  mr,
  zbl,
  jstor,
  hdl,
  pubmed,
  pmcid,
  googlebooks,
  adsbibcode,
  osti,
  ntis,
  eric,
  ssrn,
  ark,
  urn,
  lccn,
  oclc,
  libris,
  gnd,
  dnb,
  sudoc,
  sbn,
  asin,
  scopus,
  wos,
  arxivclass,
  oai,
  archiveurl,
  adsurl,
  archivedate%
}

% Handle different forms of ISSNs
\DeclareDatamodelFields[type=field, datatype=literal]{
  e-issn,
  h-issn,
  l-issn,
}
\DeclareDatamodelEntryfields[article,periodical]{
  e-issn,
  h-issn,
  l-issn,
}

% Add support for explicit placeholders
\DeclareDatamodelFields[type=field,datatype=literal]{placeholder}
\DeclareDatamodelEntryfields{placeholder}
\end{filecontents*}
